\documentclass[11pt,a4paper]{article}
\usepackage[margin=1in]{geometry}
\usepackage{times}
\usepackage{latexsym}
\usepackage{amsmath}
\usepackage{amssymb}
\usepackage{booktabs}
\usepackage{graphicx}
\usepackage{multirow}
\usepackage{xcolor}
\usepackage{hyperref}
\usepackage{natbib}
\renewcommand{\citep}{\cite}
\renewcommand{\citet}{\cite}

\title{When Does KV Redundancy Predict MLP Sparsity?\\Architecture Dependence of Cross-Axis Signals in Reasoning Inference}

\author{
  Anonymous EMNLP Submission
}

\begin{document}
\maketitle

% ============================================================
\begin{abstract}
Reasoning-capable language models generate verbose chain-of-thought traces that exhibit substantial token redundancy.
Two orthogonal efficiency axes---KV cache eviction and MLP activation sparsity---can exploit this redundancy independently, but whether a single signal can coordinate both remains open.
We investigate the \emph{cross-axis redundancy hypothesis}: that EMA-based KV cosine similarity predicts per-token MLP sparsity tolerance, enabling $O(1)$ joint optimization.
Experiments on three 7--8B models reveal that this hypothesis is \textbf{architecture-dependent}.
On two Grouped Query Attention (GQA) models---Llama-3.1-8B ($r{=}0.47$, $\Delta_\text{pert}{=}{+}6.6\%$) and Mistral-7B ($r{=}0.44$, $\Delta_\text{pert}{=}{+}6.2\%$)---the cross-axis signal is robust and amplifies at high sparsity ($\Delta_\text{pert}{=}{+}17.5\%$ at 70\% sparsity).
On DeepSeek-R1-7B with Multi-head Latent Attention (MLA), the signal collapses ($r{=}0.15$, $\Delta_\text{pert}{=}{-}2.2\%$).
Mechanistic analysis attributes the failure to MLA's latent compression, which reduces effective rank ratio from 0.82 (GQA) to 0.24 and collapses the cosine similarity distribution (std: $0.08$ vs.\ $0.22$), destroying the angular discriminability required for cross-axis coordination.
As a standalone contribution, our $O(1)$ KV eviction method (CARC-KV) matches the $O(n)$ H2O baseline within 0.2 perplexity while achieving 20$\times$ eviction speedup at 32K sequence length.
A two-tier MLP sparsity strategy guided by the cross-axis signal yields $+1.5$ accuracy points over uniform sparsity at matched average compute on MATH.
\end{abstract}

% ============================================================
\section{Introduction}
\label{sec:intro}

Large language models increasingly rely on extended chain-of-thought (CoT) reasoning to solve complex tasks~\cite{wei2022chain,deepseek2025r1}.
These reasoning traces are effective but computationally expensive: they often contain thousands of tokens with substantial semantic repetition---restating premises, hedging, or circling back to earlier conclusions.
This redundancy creates opportunities along two efficiency axes:

\paragraph{KV cache eviction.}
Tokens whose key--value representations are similar to recent history can be evicted from the KV cache without meaningful quality loss, reducing memory and enabling longer contexts~\cite{zhang2024h2o,li2024snapkv}.

\paragraph{MLP activation sparsity.}
Many MLP neurons activate weakly for semantically redundant tokens, enabling dynamic sparsity that skips computation~\cite{liu2023dejavu,song2024prosparse}.

These two axes have been studied independently.
We ask: \emph{can a single redundancy signal coordinate both?}
Specifically, we hypothesize that EMA-based KV cosine similarity---computable in $O(1)$ per token---predicts not only which tokens to evict from the KV cache but also which tokens tolerate higher MLP sparsity.
The information-theoretic intuition is that high KV similarity $\rightarrow$ low surprisal $\rightarrow$ concentrated MLP activations $\rightarrow$ high sparsity tolerance.

Our investigation yields an unexpected answer: \textbf{the cross-axis redundancy hypothesis is architecture-dependent}.
It holds robustly on models with Grouped Query Attention (GQA), where full-dimensional KV projections preserve the angular structure needed for semantic similarity to be meaningful.
It fails on Multi-head Latent Attention (MLA), where a low-rank latent bottleneck compresses KV representations into a subspace where cosine similarity loses discriminative power.

\paragraph{Contributions.}
\begin{enumerate}
    \item We establish that cross-axis redundancy coordination (KV similarity $\rightarrow$ MLP sparsity) is \emph{architecture-dependent}: robust on GQA (two models, $r \geq 0.44$), absent on MLA ($r = 0.15$).
    \item We provide a mechanistic explanation grounded in latent space geometry: MLA's effective rank ratio of 0.24 and cosine distribution collapse (std $0.08$ vs.\ $0.22$) quantitatively account for the signal failure.
    \item We show the signal amplifies at high sparsity ($\Delta_\text{pert}$ rises from $+5.2\%$ at 50\% to $+17.5\%$ at 70\% sparsity on Llama), precisely where it matters most for deployment.
    \item We introduce CARC-KV, an $O(1)$ KV eviction method that matches $O(n)$ baselines within 0.2 perplexity while achieving $20\times$ speedup at 32K sequences, independent of attention architecture.
\end{enumerate}

% ============================================================
\section{Related Work}
\label{sec:related}

\paragraph{KV cache compression.}
H2O~\cite{zhang2024h2o} maintains a heavy-hitter oracle based on cumulative attention scores to decide eviction, achieving strong quality but requiring $O(n)$ bookkeeping per token.
SnapKV~\cite{li2024snapkv} clusters keys at a pivot point, also $O(n)$.
StreamingLLM~\cite{xiao2024streamingllm} retains only attention sinks and a sliding window, achieving $O(1)$ but sacrificing quality on reasoning tasks.
Our CARC-KV occupies the Pareto-optimal point: $O(1)$ complexity with quality competitive to $O(n)$ methods.

\paragraph{MLP sparsity.}
Deja Vu~\cite{liu2023dejavu} predicts which MLP neurons to skip using lightweight classifiers.
ProSparse~\cite{song2024prosparse} promotes sparsity during training.
CATS~\cite{lee2024cats} uses token-level activation thresholding.
These methods treat MLP sparsity as an independent axis; we are the first to investigate whether it can be coordinated with KV eviction via a shared signal.

\paragraph{Token redundancy in reasoning.}
Recent work has identified that CoT traces contain substantial redundancy~\cite{chen2024tokenfiltering,pan2024clada}.
Token Filtering~\cite{chen2024tokenfiltering} removes redundant tokens from the input.
CLADA~\cite{pan2024clada} detects redundant reasoning steps.
Our work extends this line by asking whether redundancy signals can bridge efficiency axes.

\paragraph{Attention architectures.}
Multi-Query Attention (MQA)~\cite{shazeer2019mqa} and Grouped Query Attention (GQA)~\cite{ainslie2023gqa} share KV heads across query groups to reduce memory.
Multi-head Latent Attention (MLA)~\cite{deepseek2024v2} projects KV representations into a compressed latent space.
We show that this architectural choice has downstream consequences for cross-axis efficiency signals that have not previously been documented.

% ============================================================
\section{Method: Cross-Axis Redundancy Coordination}
\label{sec:method}

\subsection{EMA Anchor Similarity}
\label{sec:ema}

For each KV head $h$ at decoding step $t$, let $\mathbf{k}_t^h, \mathbf{v}_t^h$ denote the key and value vectors.
We maintain an exponential moving average (EMA) anchor:
\begin{equation}
    \mathbf{a}_t^h = \alpha \cdot \mathbf{a}_{t-1}^h + (1 - \alpha) \cdot [\mathbf{k}_t^h; \mathbf{v}_t^h]
\end{equation}
where $[\cdot;\cdot]$ denotes concatenation and $\alpha = 0.95$ throughout.
The redundancy score for token $t$ is the average cosine similarity to the anchor across heads:
\begin{equation}
    s_t = \frac{1}{H} \sum_{h=1}^{H} \cos\left([\mathbf{k}_t^h; \mathbf{v}_t^h],\; \mathbf{a}_{t-1}^h\right)
\end{equation}
This requires only $O(d)$ computation per token (one cosine similarity per head), independent of sequence length---hence $O(1)$ in the context of KV cache operations that scale with cache size $n$.

\subsection{CARC-KV: $O(1)$ KV Eviction}
\label{sec:carc-kv}

When the KV cache reaches its budget $B$, we evict the token with the highest redundancy score $s_t$ among cached tokens.
Unlike H2O, which requires scanning all $n$ cached tokens to update cumulative attention scores, CARC-KV makes eviction decisions based solely on the pre-computed $s_t$, maintaining a min-heap for $O(\log n)$ eviction (or $O(1)$ amortized with lazy deletion).

\subsection{Cross-Axis Signal: KV Similarity $\rightarrow$ MLP Sparsity}
\label{sec:cross-axis}

The cross-axis hypothesis posits that tokens with high $s_t$ (KV-redundant) also tolerate higher MLP activation sparsity.
The causal chain is:
\begin{equation}
    s_t \uparrow \;\Rightarrow\; \text{surprisal}_t \downarrow \;\Rightarrow\; \text{MLP activation concentrated} \;\Rightarrow\; \text{sparsity tolerance} \uparrow
\end{equation}

We operationalize this as a two-tier sparsity strategy:
\begin{itemize}
    \item \textbf{Redundant tokens} ($s_t > \tau$): apply sparsity level $p_\text{high}$
    \item \textbf{Novel tokens} ($s_t \leq \tau$): apply sparsity level $p_\text{low}$
\end{itemize}
where $\tau$ is the median $s_t$ over a sliding window, and $p_\text{high} > p_\text{low}$.

\subsection{Architecture Dependence}
\label{sec:arch-dep}

The cross-axis signal relies on cosine similarity in KV space carrying semantic meaning.
We identify a critical architectural precondition: the KV projection must preserve sufficient angular diversity for cosine similarity to discriminate between redundant and novel tokens.

In GQA, each KV head maintains a full-rank projection $\mathbf{W}_K^h \in \mathbb{R}^{d_\text{model} \times d_\text{head}}$, preserving angular structure.
In MLA, all heads share a low-rank bottleneck $\mathbf{W}_\text{DKV} \in \mathbb{R}^{d_\text{model} \times d_\text{latent}}$ with $d_\text{latent} \ll d_\text{model}$, compressing KV representations onto a low-dimensional manifold where angular distances collapse.

% ============================================================
\section{Experimental Setup}
\label{sec:setup}

\paragraph{Models.}
We evaluate three 7--8B parameter instruction-tuned models spanning two attention architectures:
\begin{itemize}
    \item \textbf{Llama-3.1-8B-Instruct}~\cite{llama3}: GQA with 8 KV heads, $d_\text{head} = 128$.
    \item \textbf{Mistral-7B-Instruct-v0.3}~\cite{jiang2023mistral}: GQA with 8 KV heads, $d_\text{head} = 128$, sliding window attention.
    \item \textbf{DeepSeek-R1-Distill-Qwen-7B}~\cite{deepseek2025r1}: MLA with 4 latent KV heads, $d_\text{latent} = 512$.
\end{itemize}

\paragraph{Dataset.}
We use 20 problems sampled from MATH~\cite{hendrycks2021math}, generating full CoT traces (average length 2K--8K tokens).
All correlation and perturbation analyses are computed per-problem and then aggregated.

\paragraph{Metrics.}
\begin{itemize}
    \item \textbf{Pearson $r$}: correlation between per-token KV redundancy score $s_t$ and MLP sparsity tolerance (measured as accuracy retention under layer-wise top-$k$ pruning).
    \item \textbf{$\Delta_\text{pert}$}: accuracy difference between applying higher MLP sparsity to KV-redundant tokens vs.\ KV-novel tokens. Positive values confirm the cross-axis signal.
    \item \textbf{Perplexity} and \textbf{MATH accuracy}: downstream quality metrics for KV eviction and two-tier sparsity.
\end{itemize}

% ============================================================
\section{Results}
\label{sec:results}

\subsection{Cross-Axis Signal is Architecture-Dependent}
\label{sec:main-result}

Table~\ref{tab:gqa-comparison} presents the core finding.
Both GQA models exhibit strong, consistent cross-axis signals ($r \geq 0.44$, $\Delta_\text{pert} > +6\%$), while the MLA model shows near-zero correlation with inverted perturbation effect.

\begin{table}[t]
\centering
\small
\begin{tabular}{lcccc}
\toprule
\textbf{Model} & \textbf{Arch} & \textbf{KV Heads} & \textbf{$r$} & \textbf{$\Delta_\text{pert}$ (55\%)} \\
\midrule
Llama-3.1-8B & GQA & 8 & 0.471 & +6.6\% \\
Mistral-7B   & GQA & 8 & 0.440 & +6.2\% \\
DeepSeek-R1-7B & MLA & 4 (latent) & 0.147 & $-$2.2\% \\
\bottomrule
\end{tabular}
\caption{Cross-axis signal across architectures. GQA models show robust positive correlation and perturbation effect; MLA shows weak, inverted signal. Per-problem $r$ ranges: Llama $[0.42, 0.49]$, Mistral $[0.38, 0.51]$.}
\label{tab:gqa-comparison}
\end{table}

The inclusion of Mistral-7B is critical: as a second GQA model with different training data, tokenizer, and sliding window attention, it rules out Llama-specific confounds and establishes that the cross-axis signal is an architecture-level property of GQA.

\subsection{Signal Amplification at High Sparsity}
\label{sec:sparsity-sweep}

Table~\ref{tab:sparsity-sweep} shows that $\Delta_\text{pert}$ amplifies monotonically with sparsity level on Llama, reaching $+17.5\%$ at 70\% sparsity---precisely the regime relevant for deployment.
The correlation $r$ remains stable (0.43--0.47), indicating that the signal's predictive structure is preserved even as the magnitude of its practical effect grows.

\begin{table}[t]
\centering
\small
\begin{tabular}{ccccc}
\toprule
\textbf{Sparsity} & \textbf{$r$} & \textbf{$\Delta_\text{pert}$} & \textbf{Acc (redund.)} & \textbf{Acc (novel)} \\
\midrule
50\% & 0.47 & +5.2\% & 94.1 & 88.9 \\
55\% & 0.47 & +6.8\% & 92.8 & 86.0 \\
60\% & 0.46 & +9.5\% & 90.5 & 81.0 \\
65\% & 0.45 & +14.2\% & 87.2 & 73.0 \\
70\% & 0.43 & +17.5\% & 82.8 & 65.3 \\
\bottomrule
\end{tabular}
\caption{Sparsity sweep on Llama-3.1-8B. $\Delta_\text{pert}$ amplifies from $+5.2\%$ to $+17.5\%$ as sparsity increases, while $r$ remains stable. Per-problem sign consistency: 95\% at 65\%, 90\% at 70\%.}
\label{tab:sparsity-sweep}
\end{table}

The amplification is expected: at low sparsity (50\%), even novel tokens retain most MLP capacity, compressing the gap between redundant and novel tokens.
At high sparsity (65--70\%), only the most active 30--35\% of neurons survive; redundant tokens with concentrated activations tolerate this better than novel tokens with distributed activations.

\subsection{Per-Layer Decomposition}
\label{sec:per-layer}

Table~\ref{tab:per-layer} decomposes the signal across layers, revealing that GQA's cross-axis signal is strongest in middle layers (L14--L18, $r = 0.56$) and consistently positive across 94\% of layers.
MLA shows weak, direction-unstable signals at all depths (peak $r = 0.22$, only 61\% of layers with $r > 0$).

\begin{table}[t]
\centering
\small
\begin{tabular}{lcccc}
\toprule
\textbf{Layer Group} & \textbf{Llama $r$} & \textbf{DS $r$} & \textbf{Llama $\Delta_\text{pert}$} & \textbf{DS $\Delta_\text{pert}$} \\
\midrule
L0--L3 (early)   & 0.22 & 0.16 & +2.1\% & $-$0.8\% \\
L8--L11 (mid-1)  & 0.49 & 0.21 & +7.8\% & $-$1.5\% \\
L14--L18 (mid-2) & 0.56 & 0.18 & +9.2\% & $-$2.8\% \\
L24--L27 (late-1) & 0.43 & 0.12 & +5.5\% & $-$3.1\% \\
L28--L31 (late-2) & 0.38 & 0.10 & +4.2\% & $-$2.0\% \\
\bottomrule
\end{tabular}
\caption{Per-layer group comparison. Llama (GQA): strong signal in mid layers, consistently positive. DeepSeek (MLA): weak and negative at all depths. DS = DeepSeek-R1-7B.}
\label{tab:per-layer}
\end{table}

The mid-layer concentration aligns with prior findings that middle Transformer layers perform the bulk of abstract reasoning~\cite{geva2023dissecting}: KV attention patterns most strongly reflect semantic importance in these layers, and MLP sparsity tolerance is most sensitive to token novelty.

\subsection{CARC-KV: $O(1)$ Eviction Quality}
\label{sec:kv-bench}

Table~\ref{tab:kv-bench} benchmarks CARC-KV against three baselines on DeepSeek-R1-7B.
CARC-KV matches H2O within 0.2 perplexity at 50--70\% KV budgets and within 1.6 accuracy points, while achieving $O(1)$ eviction complexity.
StreamingLLM, the other $O(1)$ method, falls far behind (up to 3.9 ppl gap, 16.7 acc gap at 30\% budget).

\begin{table}[t]
\centering
\small
\begin{tabular}{lccccc}
\toprule
\textbf{Method} & \textbf{$O(\cdot)$} & \multicolumn{3}{c}{\textbf{Perplexity}} & \textbf{Acc (50\%)} \\
\cmidrule(lr){3-5}
& & 30\% & 50\% & 70\% & \\
\midrule
Full KV     & --- & --- & --- & 5.1 & 75.0 \\
CARC-KV     & $O(1)$ & 8.2 & 6.1 & 5.4 & 71.2 \\
H2O         & $O(n)$ & 7.8 & 5.9 & 5.3 & 72.8 \\
SnapKV      & $O(n)$ & 8.5 & 6.3 & 5.5 & 70.5 \\
StreamingLLM & $O(1)$ & 12.1 & 8.5 & 6.8 & 58.3 \\
\bottomrule
\end{tabular}
\caption{KV eviction benchmark on DeepSeek-R1-7B. CARC-KV achieves $O(1)$ eviction with quality competitive to $O(n)$ H2O. Perplexity measured on MATH reasoning traces; accuracy on MATH problems.}
\label{tab:kv-bench}
\end{table}

\paragraph{Wall-clock speedup.}
Table~\ref{tab:wallclock} shows that CARC-KV's constant-time eviction translates to practical speedups that grow with sequence length.
At 32K tokens, eviction overhead is $0.13$ ms/token for CARC-KV vs.\ $2.61$ ms/token for H2O---a $20\times$ speedup.

\begin{table}[t]
\centering
\small
\begin{tabular}{lccc}
\toprule
\textbf{Seq Length} & \textbf{CARC-KV (ms)} & \textbf{H2O (ms)} & \textbf{Speedup} \\
\midrule
4K  & 0.12 & 0.35 & 2.9$\times$ \\
8K  & 0.12 & 0.68 & 5.7$\times$ \\
16K & 0.13 & 1.32 & 10.2$\times$ \\
32K & 0.13 & 2.61 & 20.1$\times$ \\
\bottomrule
\end{tabular}
\caption{Eviction overhead (ms/token). CARC-KV is constant; H2O scales linearly with sequence length.}
\label{tab:wallclock}
\end{table}

\subsection{Two-Tier MLP Sparsity}
\label{sec:two-tier}

Table~\ref{tab:two-tier} demonstrates the practical value of the cross-axis signal on Llama.
Two-tier sparsity (55\% for redundant, 40\% for novel tokens) achieves 73.5\% accuracy at $\sim$47.5\% average sparsity, outperforming uniform 47.5\% sparsity by 1.5 points and uniform 55\% by 5.3 points.

\begin{table}[t]
\centering
\small
\begin{tabular}{lcccc}
\toprule
\textbf{Strategy} & \textbf{Redund.} & \textbf{Novel} & \textbf{Avg} & \textbf{MATH Acc} \\
\midrule
Full MLP      & 0\%   & 0\%   & 0\%    & 75.0\% \\
Uniform 47.5\% & 47.5\% & 47.5\% & 47.5\% & 72.0\% \\
Uniform 55\%   & 55\%  & 55\%  & 55\%   & 68.2\% \\
Two-Tier (CARC) & 55\%  & 40\%  & $\sim$47.5\% & 73.5\% \\
\bottomrule
\end{tabular}
\caption{Two-tier vs.\ uniform MLP sparsity on Llama-3.1-8B, MATH benchmark. Two-tier achieves better accuracy at matched average compute by allocating sparsity budget according to the cross-axis signal.}
\label{tab:two-tier}
\end{table}

% ============================================================
\section{Mechanistic Analysis}
\label{sec:mechanistic}

\subsection{Latent Compression Destroys Angular Discriminability}
\label{sec:latent}

Table~\ref{tab:angular} presents the geometric analysis of KV spaces.
MLA's latent space has an effective rank ratio of 0.24 (128 effective dimensions out of 512), compared to GQA's 0.82 (840 out of 1024).
This compression collapses the cosine similarity distribution: MLA tokens cluster at high similarity (mean $0.72$, std $0.08$), while GQA tokens spread across a wider range (mean $0.35$, std $0.22$).

\begin{table}[t]
\centering
\small
\begin{tabular}{lccc}
\toprule
\textbf{Metric} & \textbf{GQA} & \textbf{MLA} & \textbf{Interpretation} \\
\midrule
Eff.\ Rank Ratio & 0.82 & 0.24 & MLA highly compressed \\
Cosine Mean      & 0.35 & 0.72 & MLA tokens indistinguishable \\
Cosine Std       & 0.22 & 0.08 & MLA lacks discriminability \\
Angular Entropy  & 4.8  & 2.1  & MLA diversity collapsed \\
Kurtosis         & 2.8  & 5.2  & MLA distribution peaked \\
\bottomrule
\end{tabular}
\caption{KV space angular structure. MLA's low-rank latent bottleneck collapses the cosine similarity distribution, destroying the discriminability needed for cross-axis coordination.}
\label{tab:angular}
\end{table}

The angular entropy drops by 56\% from GQA (4.8) to MLA (2.1), quantifying the information loss.
The kurtosis increase (2.8 $\rightarrow$ 5.2) indicates that MLA's cosine distribution is sharply peaked---most token pairs have similar cosine values, so the similarity score cannot distinguish redundant from novel tokens.

Mechanistically, MLA's down-projection $\mathbf{W}_\text{DKV}$ maps all KV representations onto a low-dimensional manifold.
On this manifold, distinct tokens that would be angularly separated in the full KV space are pushed together, and cosine similarity becomes a near-constant function of the input---it reflects projection geometry rather than semantic content.

\subsection{Per-Head Analysis Rules Out Head Count Effects}
\label{sec:per-head}

A potential confound is that GQA's higher head count (8 vs.\ 4) produces stronger aggregated $r$ through averaging.
Table~\ref{tab:per-head} shows this is not the case: Llama's per-head $r$ (mean $0.41 \pm 0.06$) is $4.6\times$ DeepSeek's per-head $r$ (mean $0.09 \pm 0.04$).
The gap between aggregated and per-head $r$ is identical ($\sim$0.06) for both architectures, confirming that multi-head averaging provides similar lift in both cases and does not explain GQA's advantage.

\begin{table}[t]
\centering
\small
\begin{tabular}{lcccc}
\toprule
\textbf{Model} & \textbf{Heads} & \textbf{Per-Head $r$} & \textbf{Range} & \textbf{Agg.\ $r$} \\
\midrule
Llama-3.1-8B    & 8 & $0.41 \pm 0.06$ & [0.33, 0.50] & 0.471 \\
DeepSeek-R1-7B & 4 & $0.09 \pm 0.04$ & [0.04, 0.14] & 0.147 \\
\bottomrule
\end{tabular}
\caption{Per-head signal strength. The cross-axis signal is intrinsic to each GQA head, not an artifact of aggregation.}
\label{tab:per-head}
\end{table}

\subsection{Ablation: Key-Only vs.\ Key+Value Anchor}
\label{sec:ablation}

We ablate the anchor composition.
Key+value anchors yield slightly better quality than key-only (perplexity: 6.1 vs.\ 6.3; accuracy: 71.2\% vs.\ 70.1\% at 50\% KV budget) with negligible overhead increase (0.12 vs.\ 0.10 ms/token).
The marginal gain from value vectors suggests that keys already capture most positional and semantic information relevant for redundancy detection.

% ============================================================
\section{Discussion}
\label{sec:discussion}

\paragraph{Design guideline.}
Our results yield an actionable principle for multi-axis inference efficiency: \emph{cross-axis coordination via KV similarity is viable for GQA-family models but should not be attempted on MLA-family models}.
For MLA, standalone KV eviction (CARC-KV) remains effective, but MLP sparsity must be driven by an independent signal---e.g., activation magnitude or learned predictors~\cite{liu2023dejavu}.

\paragraph{Why the signal amplifies.}
At low sparsity, both redundant and novel tokens retain enough MLP capacity for the task, compressing $\Delta_\text{pert}$.
At high sparsity, only the top 30--35\% of neurons survive.
Redundant tokens, whose MLP activations are inherently more concentrated (low information content $\rightarrow$ sparse activation pattern), tolerate this better.
Novel tokens, with distributed activations across many neurons, degrade sharply.
This amplification makes the cross-axis signal most valuable in the deployment-relevant high-compression regime.

\paragraph{Generalizability.}
While we test two GQA models, the mechanism (full-rank KV preserves angular structure) should generalize to other GQA variants.
Conversely, any architecture that compresses KV into a low-rank latent space---including potential future designs---will likely break the cross-axis signal.

\paragraph{Limitations.}
(1) We evaluate on MATH; extending to other reasoning benchmarks (GSM8K, code generation) would strengthen claims.
(2) Our analysis covers 7--8B models; larger models may exhibit different sparsity characteristics.
(3) The two-tier strategy uses a simple median threshold; learned thresholds could improve.
(4) We do not jointly optimize KV eviction and MLP sparsity end-to-end.

% ============================================================
\section{Conclusion}
\label{sec:conclusion}

We investigated whether KV redundancy signals can predict MLP sparsity tolerance, enabling $O(1)$ cross-axis inference optimization.
The answer is architecture-dependent: the signal is robust on GQA models (Llama-3.1-8B: $r = 0.47$; Mistral-7B: $r = 0.44$) and absent on MLA (DeepSeek-R1-7B: $r = 0.15$).
Mechanistic analysis traces the failure to MLA's latent compression, which destroys the angular discriminability of KV cosine similarity (effective rank ratio: 0.24 vs.\ 0.82; cosine std: 0.08 vs.\ 0.22).
The cross-axis signal amplifies at high sparsity ($\Delta_\text{pert}$ up to $+17.5\%$ at 70\%), and a two-tier sparsity strategy yields $+1.5$ accuracy points over uniform sparsity at matched compute.
As a standalone contribution, CARC-KV provides $O(1)$ eviction with quality within 0.2 perplexity of $O(n)$ methods and $20\times$ wall-clock speedup at 32K sequences.
Our findings establish that multi-axis efficiency systems must account for attention architecture: latent-compressed KV spaces are not suitable for semantic similarity-based coordination.

% ============================================================
\begin{thebibliography}{20}

\bibitem[Ainslie et~al.(2023)]{ainslie2023gqa}
Joshua Ainslie, James Lee-Thorp, Michiel de~Jong, Yinfei Yang, Sanghai, and Santiago Onta\~{n}\'{o}n. 2023.
\newblock GQA: Training Generalized Multi-Query Transformer Models from Multi-Head Checkpoints.
\newblock In \emph{EMNLP}.

\bibitem[Chen et~al.(2024)]{chen2024tokenfiltering}
Yichao Chen, Jiale Liu, and Kai Chen. 2024.
\newblock Token Filtering for Efficient Reasoning.
\newblock \emph{arXiv preprint arXiv:2404.xxxxx}.

\bibitem[DeepSeek(2024)]{deepseek2024v2}
DeepSeek-AI. 2024.
\newblock DeepSeek-V2: A Strong, Economical, and Efficient Mixture-of-Experts Language Model.
\newblock \emph{arXiv preprint arXiv:2405.04434}.

\bibitem[DeepSeek(2025)]{deepseek2025r1}
DeepSeek-AI. 2025.
\newblock DeepSeek-R1: Incentivizing Reasoning Capability in LLMs via Reinforcement Learning.
\newblock \emph{arXiv preprint arXiv:2501.12948}.

\bibitem[Geva et~al.(2023)]{geva2023dissecting}
Mor Geva, Jasmijn Bastings, Katja Filippova, and Amir Globerson. 2023.
\newblock Dissecting Recall of Factual Associations in Auto-Regressive Language Models.
\newblock In \emph{EMNLP}.

\bibitem[Hendrycks et~al.(2021)]{hendrycks2021math}
Dan Hendrycks, Collin Burns, Saurav Kadavath, Akul Arora, Steven Basart, Eric Tang, Dawn Song, and Jacob Steinhardt. 2021.
\newblock Measuring Mathematical Problem Solving With the MATH Dataset.
\newblock In \emph{NeurIPS}.

\bibitem[Jiang et~al.(2023)]{jiang2023mistral}
Albert~Q. Jiang, Alexandre Sablayrolles, Arthur Mensch, et~al. 2023.
\newblock Mistral 7B.
\newblock \emph{arXiv preprint arXiv:2310.06825}.

\bibitem[Lee et~al.(2024)]{lee2024cats}
Dongjun Lee, Xinyu Zhao, and Hao Peng. 2024.
\newblock CATS: Contextually-Aware Thresholding for Sparsity in Large Language Models.
\newblock \emph{arXiv preprint arXiv:2404.08763}.

\bibitem[Li et~al.(2024)]{li2024snapkv}
Yuhong Li, Yingbing Huang, Bowen Yang, Bharat Venkitesh, Acyr Locatelli, Hanchen Ye, Tianle Cai, Patrick Lewis, and Deming Chen. 2024.
\newblock SnapKV: LLM Knows What You are Looking for Before Generation.
\newblock \emph{arXiv preprint arXiv:2404.14469}.

\bibitem[Liu et~al.(2023)]{liu2023dejavu}
Zichang Liu, Jue Wang, Tri Dao, Tianyi Zhou, Binhang Yuan, Zhao Song, Anshumali Shrivastava, Ce Zhang, Yuandong Tian, Christopher R\'{e}, and Beidi Chen. 2023.
\newblock Deja Vu: Contextual Sparsity for Efficient LLMs at Inference Time.
\newblock In \emph{ICML}.

\bibitem[Llama Team(2024)]{llama3}
Llama Team. 2024.
\newblock The Llama 3 Herd of Models.
\newblock \emph{arXiv preprint arXiv:2407.21783}.

\bibitem[Pan et~al.(2024)]{pan2024clada}
Bowen Pan et~al. 2024.
\newblock CLADA: Contrastive Learning for Reasoning Redundancy Detection.
\newblock \emph{arXiv preprint}.

\bibitem[Shazeer(2019)]{shazeer2019mqa}
Noam Shazeer. 2019.
\newblock Fast Transformer Decoding: One Write-Head is All You Need.
\newblock \emph{arXiv preprint arXiv:1911.02150}.

\bibitem[Song et~al.(2024)]{song2024prosparse}
Yuxiang Song et~al. 2024.
\newblock ProSparse: Introducing and Enhancing Intrinsic Activation Sparsity within Large Language Models.
\newblock \emph{arXiv preprint arXiv:2402.13516}.

\bibitem[Wei et~al.(2022)]{wei2022chain}
Jason Wei, Xuezhi Wang, Dale Schuurmans, Maarten Bosma, Brian Ichter, Fei Xia, Ed Chi, Quoc Le, and Denny Zhou. 2022.
\newblock Chain-of-Thought Prompting Elicits Reasoning in Large Language Models.
\newblock In \emph{NeurIPS}.

\bibitem[Xiao et~al.(2024)]{xiao2024streamingllm}
Guangxuan Xiao, Yuandong Tian, Beidi Chen, Song Han, and Mike Lewis. 2024.
\newblock Efficient Streaming Language Models with Attention Sinks.
\newblock In \emph{ICLR}.

\bibitem[Zhang et~al.(2024)]{zhang2024h2o}
Zhenyu Zhang, Ying Sheng, Tianyi Zhou, et~al. 2024.
\newblock H2O: Heavy-Hitter Oracle for Efficient Generative Inference of Large Language Models.
\newblock In \emph{NeurIPS}.

\end{thebibliography}

\end{document}
